\paragraph{\trEquivariant\trBrane。}\label{sub:equivariant-branes}
设 $(M,\varpi_M)$ 是一个\trSymplectic\trManifold，配有\trAlmostComplex{}结构
$I_M$ 与二次复\trVolume{}形式 $\eta_M^2$，并带有保持 $(\varpi_M,I_M,\eta_M^2)$ 的一个%
\trInvolution{} $\iota_M$。我们想把 $\iota_M$ 在\trLagrange\trSubmanifold{}上的自然作用，
与\trPinStructure{}的改变结合起来。这会在纯形式的层面上造成一些复杂情况，我们此刻就
来处理。一个\textbf{预\trEquivariant{}\trBrane}（pre-equivariant brane）由如下数据组成：其\textbf{主分量}%
（principal component），即一个\trLagrange\trBrane{} $L^{\#}_+=(L_+,\alpha^{\#}_+,P^{\#}_+)$；
其\textbf{\trShadow}（shadow），即\trGrading{}\trLagrange\trSubmanifold{}
$(L_-=\iota_M(L_+),\,\alpha^{\#}_-=\alpha^{\#}_+\circ\iota_M)$ 上的一个\trPinStructure{}
$P^{\#}_-$；此外还有一对\trIsomorphism{} $f_+\colon P^{\#}_+\to\iota_M^*\bigl(P^{\#}_-
\otimes\lambda^{\mathrm{top}}(TL_-)\bigr)$、$f_-\colon P^{\#}_-\to\iota_M^*\bigl(P^{\#}_+
\otimes\lambda^{\mathrm{top}}(TL_+)\bigr)$，使得\trComposition{}
\begin{equation}\label{eq:equivariant-shadow-composition}
P^{\#}_+\xrightarrow{\ f_+\ }\iota_M^*\bigl(P^{\#}_-\otimes\lambda^{\mathrm{top}}(TL_-)\bigr)
\xrightarrow{\ \iota_M^*(f_-)\otimes\mathrm{Id}\ }
P^{\#}_+\otimes\lambda^{\mathrm{top}}(TL_+)\otimes\iota_M^*\lambda^{\mathrm{top}}(TL_-)
\cong P^{\#}_+,
\end{equation}
其中最后一步用到显然的\trIsomorphism{} $\iota_M^*(TL_-)\cong TL_+$，等于该%
\trPinStructure{}的\emph{非平凡}\trAutomorphism{} $-e\in\ker(Pin_n\to O_n)$。为使术语
简洁，我们有时会说“预\trEquivariant{}\trBrane{} $L^{\#}_+$”，但这总是意味着：辅助数据 $P^{\#}_-$
与 $f_-$（或 $f_+$；当然，这两个\trIsomorphism{}中的任何一个都唯一地决定另一个）的某个
具体选取已经做出。
\begin{remark}\label{thm:equivariant-brane-remark}
每个\trLagrange\trBrane{} $L^{\#}_+$ 都能以显然的方式做成预\trEquivariant{}的：取
$P^{\#}_-=\iota_M^*(P^{\#}_+\otimes\lambda^{\mathrm{top}}(TL_+))$，令 $f_-$ 为恒同，
再相应地调整 $f_+$。
\end{remark}
在预\trEquivariant{}\trBrane{}上有一个自然的\trInvolution{} $\epsilon$，它交换 $+$、$-$ 两部分；
以缩写记号书之，即 $\epsilon(L^{\#}_+)=L^{\#}_-$。这依赖于如下事实：按另一种次序的\trComposition{}
$(\iota_M^*(f_+)\otimes\mathrm{Id})\circ f_-$ 同样等于 $-e$。\textbf{\trEquivariant{}\trBrane}%
（equivariant brane）定义为这个\trInvolution{}的一个\trFixedPoint。具体地说，它由一个
不变\trLagrange\trSubmanifold{} $L$（$\iota_M(L)=L$）、一个 $\iota_M$-不变的\trGrading{}
$\alpha^{\#}$，以及一个\trPinStructure{} $P^{\#}$ 连同\trIsomorphism{} $f\colon P^{\#}\to
\iota_M^*(P^{\#}\otimes\lambda^{\mathrm{top}}(TL))$ 给出，其“平方”
\begin{equation}\label{eq:equivariant-brane-square}
P^{\#}\xrightarrow{\ f\ }\iota_M^*\bigl(P^{\#}\otimes\lambda^{\mathrm{top}}(TL)\bigr)
\xrightarrow{\ \iota_M^*(f)\otimes\mathrm{Id}\ }
P^{\#}\otimes\lambda^{\mathrm{top}}(TL)\otimes\iota_M^*\lambda^{\mathrm{top}}(TL)
\cong P^{\#}
\end{equation}
是 $-e$。给定某个 $\iota_M$-不变的\trLagrange\trSubmanifold{} $L$，我们现在来考察它%
容许\trEquivariant{}\trBrane{}结构的条件。设 $H^1(L)$ 中的\trMaslovClass{}为零，于是 $L$ 容许一个%
\trGrading{} $\alpha^{\#}$。若 $L$ \trConnected，则差 $\alpha^{\#}\circ\iota_M-\alpha^{\#}$
是一个 $\iota_M$-反不变的常值函数，从而自动为零。另一方面，若 $L$ 有两个被 $\iota_M$
交换的\trConnectedComponent，则可以把其中之一上的 $\alpha^{\#}$ 加以调整，使这个差
消失。一般情形由这两种情形推出，要点在于：要得到\emph{不变}的\trGrading，并无额外的%
\trObstruction。
\trPinStructure{}这部分的情况要稍微有趣一些。最显然的看法，是把它当作一个分步进行的
过程：先给 $L$ 配备\trPinStructure，再试着寻找\trIsomorphism{} $f\colon P^{\#}\to
\iota_M^*(P^{\#}\otimes\lambda^{\mathrm{top}}(TL))$，最后确定它的平方。相应的%
\trObstruction{}分别位于 $H^2(L;\ZZ/2)$、$H^1(L;\ZZ/2)$ 与 $H^0(L;\ZZ/2)$ 之中，但后
两个是次级\trObstruction：它们依赖于先前做出的选取，并且不取遍一切值。另一种做法，是
把所有这些方面合并成一个单一而更具\trCanonical{}性的\trObstruction，即\trEquivariant{}\trCohomology{}类
\begin{equation}\label{eq:equivariant-obstruction-class}
w^{\mathrm{eq}}_2(L)+w^{\mathrm{eq}}_1(L)u+u^2\in H^2_{\ZZ/2}(L;\ZZ/2).
\end{equation}
这里 $u$ 是 $H^*(B\ZZ/2;\ZZ/2)=\ZZ/2[u]$ 的\trGenerator；而 $w^{\mathrm{eq}}_k(L)$ 是
$TL$ 的\trEquivariant{} \trStiefelWhitney{}类，它们依赖于 $TL$ 作为\trEquivariant{}\trVectorBundle{}的结构。在
$\iota_M|L$ 没有\trFixedPoint{}的情形，有一个\trStabilization{}技巧，把我们的问题化归为关于非%
\trEquivariant{}\trPinStructure{}的标准问题，这也解释了 \eqref{eq:equivariant-obstruction-class}。
具体地说，设 $\pi\colon L\to\bar{L}$ 是到\trQuotient{} $\bar{L}=L/(\ZZ/2)$ 的\trCovering，
而 $\beta\to\bar{L}$ 是与之相伴的\trRealLineBundle。假设我们有 $T\bar{L}\oplus\beta\oplus
\beta\oplus\beta$ 上的一个\trPinStructure{} $\bar{P}^{\#}$，这意味着
\begin{equation}\label{eq:pin-quotient-criterion}
w_2(T\bar{L}\oplus\beta\oplus\beta\oplus\beta)
=w_2(T\bar{L})+w_1(T\bar{L})w_1(\beta)+w_1(\beta)^2\in H^2(\bar{L};\ZZ/2)
\end{equation}
必须为零。取\trPullback{} $\pi^*\bar{P}^{\#}$——它是 $TL\oplus\RR^3$ 上的一个\trPinStructure——
并设 $P^{\#}$ 是它在 $TL$ 上的限制。该\trPullback{}带有一个显然的\trAutomorphism，\trCovering{}
$D(\iota_M|L)\oplus(-\mathrm{Id}_{\RR^3})$（就\trFrame{}丛上的作用而言）。我们想把它同一个\trCovering{}
$\mathrm{Id}_{TL}\oplus(-\mathrm{Id}_{\RR^3})$ 的辅助\trAutomorphism{} $g$ \trComposition{}，然后把%
\trComposition{}限制到 $P^{\#}$ 上。由 \eqref{eq:aut-pin-diagram} 与 \cref{thm:pin-involution-preimage}，
$g$ 的选取对应于 $\RR^3\otimes\lambda^{\mathrm{top}}(TL\oplus\RR^3)$ 的一个\trOrientation，
而这与 $TL$ 自身的\trOrientation{}是同一回事。把这个因素考虑进来之后（比如选取局部%
\trOrientation，并把所得的 $P^{\#}$ 之局部\trAutomorphism{}拼合起来），结果便是：$P^{\#}$
带有一个自然的\trAutomorphism{} $f\colon P^{\#}\to\iota_M^*(P^{\#}\otimes\lambda^{\mathrm{top}}(TL))$。
$\pi^*\bar{P}^{\#}$ 原来的\trAutomorphism{}按定义是个\trInvolution；而 $g^2=-e$，这是
\cref{thm:pin-involution-preimage} 在 $\mathrm{rank}(F^{\mathrm{anti}})=3$ 情形的结论；
并且两者可交换，这意味着 \eqref{eq:equivariant-brane-square} 具有所需的性质。这一论证
反过来同样奏效，从而表明：\trEquivariant{}\trBrane{}结构的存在性蕴含 \eqref{eq:pin-quotient-criterion}
的消失性——在 $H^*_{\ZZ/2}(L)\cong H^*(\bar{L})$ 之下，它确实对应于
\eqref{eq:equivariant-obstruction-class}。
至于一般的（非自由）情形，可以经由\trBorel{}构造把任意的 $\ZZ/2$-作用变成自由作用；
这立即可见 \eqref{eq:equivariant-obstruction-class} 的消失性是 $L$ 上存在\trEquivariant{}\trBrane{}
结构的一个必要条件。充分性也可以按此证明，但需要对\trBorel{}构造上的\trPinStructure{}
多作一点讨论，我们在此从略。
为了补充前面的讨论，我们现在沿不变\trLagrange\trSubmanifold{} $L$ 的\trFixedPoint{}集
$L^{\mathrm{fix}}$，对情况作更直接的考察。把 $TL|L^{\mathrm{fix}}=F^{\mathrm{inv}}\oplus
F^{\mathrm{anti}}$ \trSplitting{}成 $D\iota_M$-不变与反不变两部分。假设我们在 $L$ 上有一个%
\trPinStructure{} $P^{\#}$，连同\trIsomorphism{} $f\colon P^{\#}\to\iota^*(P^{\#}\otimes
\lambda^{\mathrm{top}}(TL))$。由 \eqref{eq:aut-pin-diagram}，$f$ 限制到 $L^{\mathrm{fix}}$
上，由 $\mathrm{Pin}(TL\otimes\lambda^{\mathrm{top}}(TL))\times_{\ZZ/2}S(\lambda^{\mathrm{top}}(TL))$
的一个\trCovering{} $D(\iota_M|L)\otimes\mathrm{Id}\in\mathrm{O}(TL\otimes\lambda^{\mathrm{top}}(TL))$
的\trSection{}给出。再次诉诸 \cref{thm:pin-involution-preimage}，可以发现这类\trSection{}
对应于 $(F^{\mathrm{anti}}\otimes\lambda^{\mathrm{top}}(TL))\oplus\lambda^{\mathrm{top}}(TL)$
的\trOrientation，或者\trEquivalence{}地，对应于
\begin{equation}\label{eq:normal-fix-trivialization}
\lambda^{\mathrm{top}}(F^{\mathrm{anti}})^{\otimes\mathrm{rank}(F^{\mathrm{anti}})}
\otimes\lambda^{\mathrm{top}}(F^{\mathrm{inv}})^{\otimes\mathrm{rank}(F^{\mathrm{anti}})+1}
\longrightarrow L^{\mathrm{fix}}.
\end{equation}
的\trTrivialization{}。此外，$f$ 的平方将是 $e$，若 $\frac{1}{2}\,\mathrm{rank}(F^{\mathrm{anti}})
(\mathrm{rank}(F^{\mathrm{anti}})+1)$ 为偶数；而为 $-e$，若它为奇数。这里的平方按
\eqref{eq:equivariant-brane-square} 的意义理解，这正是与 \eqref{eq:pin-preimage-square}
相比会多出一个正负号 $(-1)^{\mathrm{rank}(F^{\mathrm{anti}})}$ 的原因；它经由\trPairing{}
$\iota_M^*\lambda^{\mathrm{top}}(TL)\otimes\lambda^{\mathrm{top}}(TL)\to\RR$ 的出现而来自
$\det(D(\iota_M|L))$。顺带指出，这个论证的大部分同样可以从抽象\trObstruction{}理论
\eqref{eq:equivariant-obstruction-class} 出发来进行，只要注意到
\begin{equation*}
\begin{aligned}
\bigl(w^{\mathrm{eq}}_2(L)+w^{\mathrm{eq}}_1(L)u+u^2\bigr)\big|L^{\mathrm{fix}}
&=w_2(F)+\bigl(\mathrm{rank}(F^{\mathrm{anti}})w_1(F^{\mathrm{anti}})\\
&\quad+(\mathrm{rank}(F^{\mathrm{anti}})+1)w_1(F^{\mathrm{inv}})\bigr)u\\
&\quad+\Bigl(\tfrac{\mathrm{rank}(F^{\mathrm{anti}})(\mathrm{rank}(F^{\mathrm{anti}})+1)}{2}
+1\Bigr)u^2\in H^*(L^{\mathrm{fix}};\ZZ/2)[u].
\end{aligned}
\end{equation*}
对我们的应用而言，最终只有一个特殊情形是相关的，我们把它记录下来以备后用：
\begin{lemma}\label{thm:codim-one-iso-coorientations}
设 $L$ 是 $M$ 的一个 $\iota_M$-不变\trLagrange\trSubmanifold。设 $L$ \trConnected，且其%
\trFixedPoint{}集 $L^{\mathrm{fix}}$ 依旧\trConnected、非空，并具有\trCodimension{}$1$。此外，设我们在
$L$ 上有一个\trIsomorphic{}于 $\iota_M^*(P^{\#}\otimes\lambda^{\mathrm{top}}(TL))$ 的%
\trPinStructure{} $P^{\#}$。那么，\trIsomorphism{} $f\colon P^{\#}\to\iota_M^*
(P^{\#}\otimes\lambda^{\mathrm{top}}(TL))$ 的选取，与 $L^{\mathrm{fix}}\subset L$ 的%
\trCoorientation{}\trBijective{}对应。此外，任何这类\trIsomorphism{}的平方（按
\eqref{eq:equivariant-brane-square} 的意义取）是 $-e$。
\noqed
\end{lemma}
显然，$M$ 的 $\ZZ/2$-\trEquivariant{}\trGrading{}辛\trAutomorphism{}作用在预\trEquivariant{}\trBrane{}的集合上，其作用方式与%
\trInvolution{} $\epsilon$ 严格交换，特别是保持\trEquivariant{}\trBrane。也可以用\trEquivariant{}\trLineBundle{}
$\beta\to M$ 来作用：给定这样一个\trLineBundle{}和一个预\trEquivariant{}\trBrane，把 $P^{\#}_\pm$
换成 $P^{\#}_\pm\otimes\beta|L_\pm$，并利用由 $\beta$ 的\trEquivariant{}结构给出的\trIsomorphism{}
$\beta|L_\pm\cong\iota^*(\beta|L_\mp)$，把 $f_\pm$ 变成新的一对\trPinStructure{}之间的%
\trIsomorphism。比如，取 $\beta=\RR\times M$ 带有非平凡的 $\ZZ/2$-作用（对应于 $\ZZ/2$
的非平凡实\trCharacter）：与之作张量，相当于把 $f_\pm$ 都乘上 $-e$。回到一般的 $\beta$，
把这个作用施于一个\trEquivariant{}\trBrane，并查看\trFixedPoint{}集 $L^{\mathrm{fix}}\subset L$。其效果是：
用 $\ZZ/2$ 在 $\beta|L^{\mathrm{fix}}$ 上作用的正负号，去改变 \eqref{eq:normal-fix-trivialization}
的\trTrivialization{}。
